\pdfoutput=1
\ifdefined\pdfinfoomitdate\pdfinfoomitdate=1\fi
\ifdefined\pdftrailerid\pdftrailerid{}\fi
\documentclass[11pt]{article}
\usepackage[T1]{fontenc}
\usepackage[utf8]{inputenc}
\usepackage{lmodern}
\usepackage{amsmath,amssymb,amsthm}
\usepackage[margin=1in]{geometry}
\usepackage{booktabs,array}
\usepackage{xurl}
\usepackage[unicode,colorlinks=true,linkcolor=blue,citecolor=blue,urlcolor=blue]{hyperref}
\newtheorem{theorem}{Theorem}[section]
\newtheorem{lemma}[theorem]{Lemma}
\newtheorem{proposition}[theorem]{Proposition}
\newtheorem{corollary}[theorem]{Corollary}
\theoremstyle{definition}
\newtheorem{definition}[theorem]{Definition}
\theoremstyle{remark}
\newtheorem{remark}[theorem]{Remark}
\numberwithin{equation}{section}
\setlength{\parskip}{0.25em}
\setlength{\emergencystretch}{2em}
\allowdisplaybreaks
\title{Strongly Exponential Sensitive Monotone Bounds in VP\\
via Eulerian Conditioning}
\author{Ying Xie\\Kennesaw State University\\
\texttt{yxie2@kennesaw.edu}}
\date{September 15, 2026}
\hypersetup{pdftitle={Strongly Exponential Sensitive Monotone Bounds in VP via Eulerian Conditioning},
pdfauthor={Ying Xie},
pdfsubject={Monotone arithmetic complexity},
pdfcreator={},pdfproducer={}}
\begin{document}
\maketitle
\begin{abstract}
We present a strongly exponential sensitive monotone circuit lower bound
for an explicit family of arborescence polynomials in uniform VP.
If \(N\) is the actual number of variables, \(P_H\) is the arborescence
polynomial, and \(F_H\) is the product of the outgoing-choice group sums,
then \(\operatorname{monSIZE}(F_H\pm\epsilon P_H)\ge 2^{bN}\)
for fixed \(a,b>0\) and \(2^{-aN}\le\epsilon\le1/2\).
The circuits may use arbitrary nonnegative real constants and unrestricted
fan-out. The construction conditions independent local permutations of
an expander's oriented edges to form one Eulerian cycle.
The BEST theorem and two spectral trace identities bound the conditioning
cost, while permutation exposure supplies a concentration estimate strong
enough to absorb that cost. A fixed clique blow-up realizes an
inner-product gadget and yields one distribution with exponentially small
discrepancy for every balanced partition of the variable groups.
The resulting coefficient argument charges shared gates.
As consequences, a uniform VP family with coefficients one and two has a
sharp multiplicative approximation threshold of \(\sqrt2\), even for
approximation on the whole cube without a degree restriction, and
strongly exponential exact monotone counting complexity on Boolean inputs.
\end{abstract}

\section{Introduction}\label{sec:intro}

Monotone arithmetic circuits compute polynomials using addition,
multiplication, and nonnegative constants. General arithmetic circuits
can additionally exploit cancellation. Sensitive lower bounds ask whether
a small coefficient perturbation of an easy polynomial can already force
a large monotone circuit. The perturbation parameter is part of the
statement: hardness at a fixed perturbation does not by itself imply
hardness for exponentially small perturbations.

Hrube\v{s} introduced a systematic study of this phenomenon
\cite{Hrubes2020}. Chattopadhyay, Datta, and Mukhopadhyay related
sensitive monotone lower bounds to communication complexity
\cite{CDM2021}. Chattopadhyay, Datta, Ghosal, and Mukhopadhyay
\cite{CDGM2022} proved ordinary \(2^{\Omega(N)}\) monotone lower bounds
for sparse spanning-tree polynomials in VP, and sensitive
\(2^{\Omega(\sqrt N)}\) bounds for a complete-graph family in VP.
Here and throughout this comparison, \(N\) counts actual input variables.
Their Section~7 asks for strongly exponential sensitive bounds.
Chattopadhyay, Ghosal, and Mukhopadhyay subsequently obtained such
sensitive bounds for a graph-inner-product family in VNP
\cite{CGM2022}; that result does not supply a VP upper bound.

\subsection{Main result}

For a graph \(H\) with a distinguished root, assign one variable
\(x_{z,w}\) to each arc from a nonroot vertex \(z\) to a neighbor \(w\).
Let \(F_H\) be the product, over nonroot vertices, of the sums of
their outgoing variables. Let \(P_H\) sum precisely those full-group
monomials whose chosen arcs lead every vertex to the root.
The matrix-tree determinant computes \(P_H\) by polynomial-size
general arithmetic circuits.

\begin{theorem}[Asymptotic form]\label{thm:asymptotic}
There are positive constants \(a,b\) and explicit polynomial families
\((F_N,P_N)\) on exactly \(N\) variables, for every sufficiently large \(N\),
such that \(F_N\) has an \(O(N)\)-size monotone circuit, \(P_N\) is in
uniform VP, and
\[
\operatorname{monSIZE}(F_N\pm\epsilon P_N)\ge 2^{bN}
\qquad\text{whenever}\qquad 2^{-aN}\le\epsilon\le\tfrac12.
\]
Each target has the full support of \(F_N\) and coefficients
\(1\) and \(1\pm\epsilon\). The lower bound allows arbitrary
nonnegative real constants and unrestricted sharing.
\end{theorem}

The core family is given by a fixed clique blow-up of a simple
constant-degree expander. Theorem~\ref{thm:finite} provides numerical
constants and a finite cutoff. Section~\ref{sec:explicit} supplies a
uniform explicit family and the padding argument that proves
Theorem~\ref{thm:asymptotic}. The constants are conservative; degree
and blow-up factors are fixed before the family size tends to infinity.

The point is the simultaneous combination of VP membership,
strongly exponential size in actual variables, and sensitivity down
to an exponentially small perturbation. The ordinary strongly
exponential VP separation is already known \cite{CDGM2022}.
No claim about Boolean circuit lower bounds or about separating
general arithmetic complexity classes follows from the present theorem.

\subsection{Proof mechanism}

The proof uses the discrepancy-to-coefficient framework of
\cite{CDM2021,CDGM2022}. Its central construction is a distribution
on rooted neighbor maps of a sparse host which has exponentially small
rectangle discrepancy under every balanced partition of the outgoing
variable groups. A single distribution must work for all partitions;
the sampled layout itself need not be simultaneously good for all of them.

Let \(B\) be a simple \(d\)-regular base with \(v=d|V(B)|\) oriented
edge states. Independent local permutations define a cycle cover of
these states. Condition the cover to consist of one cycle.
The BEST theorem \cite{FarrellLevine2015} expresses the probability
of this event through the spanning-tree count of \(B\). The identities
for the first and second spectral traces show that the reciprocal
conditioning probability is at most
\[
v\exp\bigl(O(v/d^2)\bigr).
\]
Keeping the signed first trace, rather than replacing it by an
absolute-value estimate, is important for this bound.

A fixed number of clones per state realizes a two-track parity gadget
along this Eulerian cycle using only edges of the host. For each fixed
balanced partition, expansion guarantees a positive expected density
of stages whose first bit is held entirely on one side and whose second
bit is held entirely on the other. Revealing permutation images
individually gives a tail bound \(\exp(-\Omega(v))\). At a sufficiently
large fixed degree this tail absorbs the conditioning cost.
The retained stages give a large inner product, whose elementary
discrepancy estimate yields the desired universal bound.

Finally, a nonnegative full-group output forces productive gates
to be set-multilinear. Removing one balanced gate at all its uses
charges at most one product to that gate. This gives the circuit-size
bound even with unrestricted fan-out.

\subsection{Consequences and related models}

The polynomial \(R_H=2F_H-P_H\) has coefficients one and two on every
full-group monomial. Section~\ref{sec:approx} proves an exact
approximation threshold: any fixed symmetric multiplicative factor
below \(\sqrt2\) requires \(2^{\Omega(N)}\) monotone size, whereas
\(\sqrt2 F_H\) attains factor \(\sqrt2\) with linear size.
This holds coefficientwise and uniformly on \([0,1]^N\), without a
degree bound on the competing circuit. The passage from the cube to
coefficients uses least-homogeneous-component extraction, a standard
monotone tool \cite{JerrumSnir1982,Jukna2016}.

Section~\ref{sec:counting} extends the size bound to exact arithmetic
counting on the whole Boolean cube. We give a full-group normalization
with explicit cost \((2g+1)s+g\), where \(g\) is the number of groups.
Removing repeated consistent indicators has classical precedents
in sum-product networks \cite{Peharz2015,MartensMedabalimi2015};
the normalization principle is not claimed as a new contribution.
The full-cube equality hypothesis is essential to this application.

The main argument occupies Sections~\ref{sec:model}--\ref{sec:explicit}.
Appendix~\ref{app:best} gives the direct BEST counting argument,
Appendix~\ref{app:thin} reduces the host degree, and
Appendix~\ref{app:general} extends the construction to any fixed-degree
simple expander family with a fixed positive spectral gap after a
suitable fixed clique blow-up.



\section{Polynomials and circuit model}\label{sec:model}

All arithmetic circuits are directed acyclic graphs with binary \(+\)
and \(\times\) gates. Size counts arithmetic gates, not input occurrences.
Monotone circuits use nonnegative real constants, and fan-out is unrestricted.
General circuits may use negative rational constants. A uniform VP family
has polynomially many variables, polynomial degree, and a polynomial-time
constructible family of polynomial-size general arithmetic circuits.
When a real \(\epsilon\) is not specified by a uniform rational sequence,
it is treated as a supplied scalar parameter.

A polynomial is \emph{full-group set-multilinear} for a partition
\(X_1,\ldots,X_g\) when each monomial uses exactly one variable from each
group, with exponent one. The notation \([m]Q\) denotes a coefficient.
For a graph \(B\), the clique blow-up \(B[K_t]\) replaces each vertex
by a \(t\)-clique and each edge by a complete bipartite graph between
the corresponding cliques. All logarithms in probability estimates
are natural unless their base is displayed.

Fix integers \(d\ge734\), \(L\ge32\). Let \(B\) be a connected
simple \(d\)-regular graph on \(n\) vertices with
\begin{equation}\label{eq:1}
\lambda_2(A_B/d)\le\rho:=2/27.
\end{equation}
Write \(v=dn\). Let \(H=B[K_{dL}]\), a fixed clique blow-up
of the base graph. Its degree is \(D_H=(d+1)dL-1\).
Choose any fixed root \(r\). Give every directed arc out of a
nonroot vertex a different variable, so the actual input count is
\begin{equation}\label{eq:2}
N=(Lv-1)D_H .
\end{equation}
Define
\[
 F_H=\prod_{z\ne r}\sum_{w\in N_H(z)}x_{z,w},\qquad
 P_H=\sum_{\nu:\ \mathrm{all\ vertices\ reach}\ r}
                              \prod_{z\ne r}x_{z,\nu(z)}.
\]
The sum is over all choices of one neighbor at each nonroot vertex.
The rooted out-Laplacian determinant equals \(P_H\), so \(F_H,P_H\)
have uniform polynomial-size general arithmetic circuits when \(B\)
is explicit.

\begin{theorem}[Finite sensitive bound]\label{thm:finite} Put \(K=LD_H\), \(c=1/7680000\),
\(a=c/(2K)\), and \(b=c/(3K)\). For every \(B\) satisfying \eqref{eq:1}
with \(v\ge2^{32}\),
\begin{equation}\label{eq:3}
\operatorname{monSIZE}(F_H\pm\epsilon P_H)\ge2^{bN}
                        \qquad(2^{-aN}\le\epsilon\le1/2).
\end{equation}
The constants \(a,b\) depend on the fixed \(d,L\).
The exponent is not uniform in growing degree or clone count.
\end{theorem}

The competing monotone circuits have binary addition and multiplication,
arbitrary nonnegative real constants, and unrestricted fan-out. Size
counts arithmetic gates. Their intermediate gates are not assumed
set-multilinear. The output coefficients are all positive: they are
1 and \(1\pm\epsilon\), and their support is that of the easy product
\(F_H\).

The communication domain is the full product
\(\mathcal X_H=\prod_{z\ne r}N_H(z)\). A partition assigns whole
outgoing-choice groups to the two players and is balanced if each
side contains between \((Lv-1)/3\) and \(2(Lv-1)/3\) groups.
For a distribution \(\Delta\), discrepancy is the largest absolute
signed mass
\[
     \left|\sum_{\nu\in\mathcal U\times\mathcal V}
                       \Delta(\nu)(2g_H(\nu)-1)\right|
\]
over rectangles for that partition, where \(g_H\) is the arborescence
indicator. There is no promise on allowable neighbor maps.

\subsection{Determinant upper bound}\label{sub:determinant}

Let \(h=|V(H)|\). Omit the row and column of the root and use the
row vector
\[
       \sum_{w\in N_H(z)}x_{z,w}(e_z-e_w),\qquad e_r:=0
\]
at each \(z\ne r\). Row multilinearity of the determinant shows
that the coefficient of \(m_\nu=\prod_{z\ne r}x_{z,\nu(z)}\)
is the determinant with rows \(e_z-e_{\nu(z)}\).
If every vertex reaches the root, a common ordering of rows and
columns by an acyclic order makes this matrix triangular with
diagonal one. Its determinant is one.
Otherwise choose a nonroot directed cycle and put a vector equal
to one on that cycle's entire basin and zero elsewhere. Its
difference across every selected outgoing arc is zero. It is a
nonzero kernel vector, so the determinant is zero.
Thus the determinant is exactly \(P_H\).

The determinant has uniform polynomial-size general arithmetic
circuits over rational constants. The construction of \(H\) is
uniform as described below, and \(F_H\) is a product of group sums.
Consequently \(F_H\pm\epsilon P_H\) has a polynomial-size circuit
with \(\epsilon\) as a parameter or a supplied constant.
For an unambiguously uniform family over the rationals one may
fix \(\epsilon=1/2\), or set
\(\epsilon_N=2^{-\lceil aN/2\rceil}\), which lies in the proved
interval for sufficiently large \(N\) and has \(O(N)\)-bit
rational representation.
No polynomial-size set-multilinear circuit or polynomial-size
formula upper bound is being asserted for \(P_H\).

\begin{center}
\begin{tabular}{ll}
\toprule
Parameter & Meaning\\
\midrule
\(n\) & Number of base vertices\\
\(d\) & Fixed base degree\\
\(v=dn\) & Number of oriented edge states\\
\(L\) & Fixed number of clones per state\\
\(h=Lv\) & Number of host vertices\\
\(D_H=(d+1)dL-1\) & Host degree\\
\(N=(h-1)D_H\) & Actual variable count\\
\bottomrule
\end{tabular}
\end{center}


\section{Eulerian distribution and parity embedding}\label{sec:distribution}

\subsection{Local permutations and the conditioning probability}\label{sub:cover}
Let \(S\) be the oriented edges of \(B\), so \(|S|=v\).
Independently at each base vertex \(w\), choose a uniform permutation
\(\sigma_w\) of its neighbors. Set
\begin{equation}\label{eq:4}
\pi(u,w)=(w,\sigma_w(u)).
\end{equation}
This is a permutation of \(S\). Immediate reversals are allowed.
Let \(E\) be the event that \(\pi\) is a single cycle.

Write \(t_o(B)\) for the number of spanning trees of \(B\), oriented
toward a fixed base vertex \(o\). The BEST theorem \cite{FarrellLevine2015} gives exactly
\begin{equation}\label{eq:5}
\Pr[E]=\frac{t_o(B)((d-1)!)^n}{(d!)^n}
                          =\frac{t_o(B)}{d^n}.
\end{equation}
A single cyclic transition configuration corresponds to exactly one
Eulerian tour starting with any prescribed directed edge; hence there
is no extra factor \(d\), \(n\), or \(v\) in the numerator of \eqref{eq:5}.
A direct counting proof and an efficient sampler appear in Appendix~\ref{app:best}.

Let \(\theta_1=1,\theta_2,\ldots,\theta_n\) be the normalized
adjacency eigenvalues. The matrix-tree formula yields
\begin{equation}\label{eq:6}
\Pr[E]=\frac1v\prod_{i=2}^n(1-\theta_i).
\end{equation}
Simplicity and regularity give the exact identities
\begin{equation}\label{eq:7}
\sum_{i=2}^n\theta_i=-1,\qquad
        \sum_{i=2}^n\theta_i^2=n/d-1=v/d^2-1.
\end{equation}
For every \(0\le\rho<1\) and \(-1\le x\le\rho\),
\begin{equation}\label{eq:8}
\log(1-x)\ge -x-\frac{x^2}{2(1-\rho)}.
\end{equation}
Indeed, if
\(f(x)=\log(1-x)+x+x^2/[2(1-\rho)]\), then
\[
                 f'(x)=\frac{x(\rho-x)}{(1-\rho)(1-x)}.
\]
Thus \(f\) decreases to \(f(0)=0\) on \([-1,0]\) and increases
on \([0,\rho]\). This also covers a bipartite eigenvalue \(-1\).
Using \eqref{eq:7}--\eqref{eq:8} in \eqref{eq:6}, and \(1/[2(1-\rho)]=27/50\), gives
\begin{equation}\label{eq:9}
\Pr[E]\ge
             \frac1v\exp\left(-\frac{27v}{50d^2}\right).
\end{equation}
The linear trace term in \eqref{eq:7} is essential: replacing it by the
absolute sum of eigenvalues would give a much more expensive
conditioning bound.

\subsection{A direct parity gadget}\label{sub:gadget}
Index the states by their tail and a neighbor label,
\(S=V(B)\times[d]\). Divide each \(dL\)-clone base cluster
of \(H\) into \(d\) labeled state clusters \(I_s\) of size \(L\).
If \(s=(u,w)\) and \(\pi(s)=(w,z)\), the two tails \(u,w\)
are adjacent in \(B\). Thus every original transition has every
needed physical intercluster edge in \(H\), including reversals.
No routing graph is used.

Fix a state \(s_0\) whose cluster contains the root \(r\).
Choose four distinct uniform labeled clones
\[
                      a_{s,0},a_{s,1},b_{s,0},b_{s,1}
\]
in every state cluster; exclude \(r\) in its own cluster.
Choices in different clusters are independent. Sample \(\pi\)
from the original local-permutation distribution conditioned on
\(E\), independently of these clone choices.

For each state \(u\ne s_0\), sample independent uniform bits
\(x_u,y_u\), and write \(w=\pi(u)\). For each track
\(\ell\in\{0,1\}\), set
\begin{equation}\label{eq:10}
a_{u,\ell}\longmapsto
 \begin{cases}a_{w,\ell},&x_u=0,\\ b_{w,\ell},&x_u=1,\end{cases}
 \qquad
 b_{w,\ell}\longmapsto a_{w,\ell\oplus y_u}.
\end{equation}
The first two destinations are different labeled clones, so both
bit choices are different valid outgoing arcs.
Set \(c_0=a_{s_0,0}\), \(c_1=a_{s_0,1}\), and
\begin{equation}\label{eq:11}
c_0\longmapsto a_{\pi(s_0),0},
                         \qquad c_1\longmapsto r.
\end{equation}
The two \(b\)-vertices in \(I_{\pi(s_0)}\) point to the same-index
\(a\)-vertices there. Every remaining nonroot clone in \(I_s\)
points to \(a_{s,0}\), using its clique edge.

For every admissible layout and every assignment of the auxiliary
bits, the following deterministic identity holds:
\begin{equation}\label{eq:12}
g_H(\nu)=\sum_{u\ne s_0}x_u y_u\pmod2.
\end{equation}
\begin{proof}[Proof of the parity identity]
Fix an admissible layout and all the auxiliary bits. Since the
conditioning event \(E\) makes \(\pi\) a single cycle, enumerate the
states as
\[
 s_i=\pi^i(s_0)\quad(0\le i\le v-1),\qquad s_v=s_0.
\]
Here \(v\ge2\), as the underlying base graph is simple and has
positive degree. We follow the active stages in the order
\(s_1,\ldots,s_{v-1}\); the transition out of \(s_0\) is instead
the special connector and root exit in \eqref{eq:11}.

\emph{The outgoing map is well defined.}
Every \(a\)-vertex outside \(I_{s_0}\) receives exactly one rule
from \eqref{eq:10}, and the two \(a\)-vertices in \(I_{s_0}\) receive
the rules in \eqref{eq:11}. Because \(\pi\) is a permutation, every
\(b\)-vertex outside \(I_{s_1}\) has a unique predecessor state
and receives exactly one rule from \eqref{eq:10}. The two \(b\)-vertices
in \(I_{s_1}\) have the separately specified, fixed rules.
Every other nonroot vertex is a filler vertex. Thus each nonroot
vertex has exactly one outgoing arc, and the root has none.
The intercluster arcs follow adjacent tails in \(B\); all other
arcs lie in a cluster clique and have distinct endpoints.
In particular, the rules neither conflict nor use a loop.

\emph{One stage and the two paths.}
Put \(t_i=x_{s_i}y_{s_i}\) for \(1\le i\le v-1\), and let
\(P=t_1\oplus\cdots\oplus t_{v-1}\).
For either incoming rail \(\ell\), the path from
\(a_{s_i,\ell}\) to the next \(a\)-vertex is exactly one of
the following:
\[
\begin{array}{ccl}
x_{s_i}=0 &:&
 a_{s_i,\ell}\longmapsto a_{s_{i+1},\ell}
 \quad\text{for either value of }y_{s_i},\\[2pt]
x_{s_i}=1,\ y_{s_i}=0 &:&
 a_{s_i,\ell}\longmapsto b_{s_{i+1},\ell}
 \longmapsto a_{s_{i+1},\ell},\\[2pt]
x_{s_i}=1,\ y_{s_i}=1 &:&
 a_{s_i,\ell}\longmapsto b_{s_{i+1},\ell}
 \longmapsto a_{s_{i+1},\ell\oplus1}.
\end{array}
\]
Consequently each stage induces the permutation
\(\ell\mapsto\ell\oplus t_i\) of the two rails.
Induction on \(j\) shows that the path beginning at
\(a_{s_1,\ell}\) reaches
\(a_{s_{j+1},\,\ell\oplus t_1\oplus\cdots\oplus t_j}\)
after \(j\) active stages. In particular, the two paths,
stopped at their final \(a\)-vertex in \(I_{s_0}\), have endpoints
\[
 a_{s_1,0}\leadsto c_P,\qquad
 a_{s_1,1}\leadsto c_{1\oplus P}.
\]
These are vertex-disjoint simple paths. Indeed, the two rail
indices stay distinct at every state, the intermediate
\(b\)-vertices are distinct from all \(a\)-vertices, and the state
order \(s_1,\ldots,s_v\) contains no repetition. Together the paths
contain every \(a\)-vertex. Neither contains \(r\).

\emph{Odd parity: \(P=1\).}
The path starting on rail zero ends at \(c_1\) and exits to
\(r\). The path starting on rail one ends at \(c_0\), whose
connector leads to the start of the first path. Thus the full
forward trajectories are
\[
\begin{aligned}
 a_{s_1,0}&\leadsto c_1\longmapsto r,\\
 a_{s_1,1}&\leadsto c_0\longmapsto
 a_{s_1,0}\leadsto c_1\longmapsto r.
\end{aligned}
\]
This accounts also for the trajectory from \(c_0\): its return
to \(I_{s_1}\) does not close a cycle, since the next active
traversal ends at \(c_1\). Every \(a\)-vertex therefore reaches
the root.

\emph{Even parity: \(P=0\).}
Now the path starting on rail zero ends at \(c_0\), so the
connector closes the simple directed cycle
\[
 a_{s_1,0}\leadsto c_0\longmapsto a_{s_1,0}.
\]
This cycle does not contain \(r\). The disjoint path starting
on rail one ends at \(c_1\) and then at \(r\).
In particular, the outgoing map fails the requirement that
every nonroot vertex reach \(r\).

\emph{Vertices outside the two paths.}
Every unused \(b\)-vertex points directly to an \(a\)-vertex
in its own state cluster. This includes the two exceptional
\(b\)-vertices in \(I_{s_1}\), whose rules are fixed, and
the \(b\)-vertices bypassed at stages with \(x_{s_i}=0\).
Every filler vertex likewise points directly to an \(a\)-vertex.
Such vertices can create no additional directed cycle.
Hence in the odd case all nonroot vertices reach \(r\);
in the even case there is exactly one nonroot directed cycle.
By the definition of \(g_H\), its value is respectively
one and zero, which proves \eqref{eq:12}.

The argument holds for every admissible layout and every bit
assignment, independently of their probability. It uses no
spectral estimate, concentration bound, or finite enumeration.
\end{proof}

There are \(k=v-1\) stages. Each \(x_u\) affects only the two
outgoing choices at \(a_{u,0},a_{u,1}\), and each \(y_u\) affects
only those at \(b_{\pi(u),0},b_{\pi(u),1}\).

Let \(\Delta\) denote the pushforward of this random experiment to
the full product of neighbor maps. Thus a layout consists of the
conditioned permutation and the clone labels, and the remaining
randomness consists of the auxiliary bits \(x_u,y_u\). These auxiliary
bits are distinct from the formal arc variables \(x_{z,w}\).
The distribution is fixed before a communication partition is chosen.


\section{Universal discrepancy and the size lower bound}\label{sec:discrepancy}

\subsection{Expansion and monochromatic pairs}\label{sub:mean}
Fix any balanced partition of all nonroot physical vertices.
For analysis assign the root either color. Let \(a_s\) be the
Alice fraction in \(I_s\), and \(\mu=v^{-1}\sum_s a_s\).
For \(Lv\ge100\), \(\mu\in[33/100,67/100]\).

The one-point transition matrix for the unconditioned cover is
\(P((u,w),(w,z))=1/d\). It is doubly stochastic.
The symmetric matrix \(Q=(P+P^\top)/2\) has image in the
functions \(a(u)+b(w)\). For each base eigenvalue \(\theta\),
its matrix on the associated two lifted functions is
\[
                     \frac12
                     \begin{pmatrix}\theta&1\\1&\theta\end{pmatrix}.
\]
Dependent lifted functions are quotiented out. Every other nonzero
eigenvector is already in this image. Thus its gap above constants
is at least \((1-\rho)/2=25/54\). A possible base eigenvalue
\(-1\) contributes a lifted eigenvalue \(-1\) and causes no
problem for this upper bound on the nonconstant spectrum.
Put \(\delta_0=25/54\). With normalized inner products,
\begin{equation}\label{eq:13}
\mathbb E_{s\ {\rm uniform},\,t\sim P(s,\cdot)} a_s(1-a_t)
       =\mu-\langle a,Qa\rangle
       \ge\delta_0\mu(1-\mu).
\end{equation}
Indeed, the spectral bound gives
\(\langle a,Qa\rangle\le\mu^2+(1-\delta_0)\operatorname{Var}(a)\).
Since \(0\le a_s\le1\), we have
\(\operatorname{Var}(a)\le\mu(1-\mu)\), yielding \eqref{eq:13}.

Put \(c_L=1/L\). The probability of selecting two Alice clones
in a cluster with Alice fraction \(a=j/L\) is
\[
             p_L(a)=\frac{j(j-1)}{L(L-1)}
                    =\frac{a(a-c_L)}{1-c_L}.
\]
It satisfies \(\sqrt{p_L(a)}\ge(a-c_L)_+\).
For \(a=0\) this is immediate; for \(a\ge c_L\), square both
nonnegative sides and use
\[
 p_L(a)-(a-c_L)^2
     =\frac{c_L(a-c_L)(a+1-c_L)}{1-c_L}\ge0.
\]
For nonnegative \(A,B,c\),
\((A-c)_+(B-c)_+\ge AB-c(A+B)\): if one factor is truncated,
the right side is nonpositive; otherwise expand the left side.
Since the transition is doubly stochastic,
\(\mathbb E[a_s+1-a_t]=1\). Hence \eqref{eq:13} implies
\[
 \mathbb E\sqrt{p_L(a_s)p_L(1-a_t)}
                 \ge\delta_0\mu(1-\mu)-\frac1L
                 \ge\frac{16}{225}\qquad(L\ge32).
\]
The final value follows exactly from
\((25/54)(33\cdot67/10000)-1/32=16/225\).
Cauchy--Schwarz gives
\(\mathbb E[p_L(a_s)p_L(1-a_t)]\ge256/50625\).
Original transitions are between distinct state clusters because
the base graph has no loops, so their selected pairs are independent.

Let \(Z\) count original transitions \(u\mapsto\pi(u)\) whose
two \(a_u\)-vertices are Alice and whose two \(b_{\pi(u)}\)-vertices
are Bob. Under the unconditioned cover and clone choices,
\begin{equation}\label{eq:14}
\mathbb E Z\ge
  \frac{256}{50625}v-2
                             \ge\eta v,\qquad \eta:=1/200.
\end{equation}
The subtraction of two covers the exceptional root cluster.
The excess of the coefficient over \(1/200\) is exactly
\(23/405000\), so \(v\ge35218\) suffices for \eqref{eq:14}.

\subsection{Exposure concentration and conditioning}\label{sub:concentration}
We first work under the \emph{unconditioned} law: the local
permutations are independent and uniform, and the clone placements
are independent across state clusters and independent of the
permutations. Fix the communication partition, including the root's
analysis color. For the clone placement \(\mathcal R_s\) in \(I_s\),
write
\[
 A_s={\bf1}_{\{a_{s,0},a_{s,1}\ {\rm are\ Alice}\}},\qquad
 B_s={\bf1}_{\{b_{s,0},b_{s,1}\ {\rm are\ Bob}\}}.
\]
Then the count from \eqref{eq:14} is
\[
 Z=\sum_{s\in S} A_sB_{\pi(s)}.
\]
The auxiliary input bits are not involved in \(Z\).

\emph{Exposure order and conditional means.}
Choose deterministic orders of the state clusters, base vertices,
and neighbors of each base vertex. First reveal the \(v\) clone
placements \(\mathcal R_s\), one whole placement at a time.
Then, at each base vertex \(w\), reveal the first \(d-1\) images
of \(\sigma_w\) in the chosen neighbor order; its last image is
forced. Let \(\mathcal F_i\) be the information after step \(i\).
There are \(J=v+n(d-1)\le2v\) steps, and
\[
 Y_i=\mathbb E[Z\mid\mathcal F_i],\qquad
 Y_0=\mathbb EZ,\quad Y_J=Z
\]
is the corresponding Doob martingale.
For a fixed positive-probability history through step \(i-1\),
let \(h_i(q)\) be the conditional expectation of \(Z\) when the
next revealed value is \(q\). We show that
\(\lvert h_i(q)-h_i(q')\rvert\le2\) for all compatible \(q,q'\).

\emph{A clone-placement step.}
Compare two possible placements in \(I_s\), and couple all
unrevealed placements and all local permutations identically.
This is a coupling of the correct conditional laws by independence.
For each resulting fixed permutation \(\pi\), only the summands
indexed by \(s\) and \(\pi^{-1}(s)\) can change: these are the
unique outgoing and incoming transitions of the state.
Each summand lies in \(\{0,1\}\), so the two completed values of
\(Z\) differ by at most two. Averaging over this coupling gives
the required bound for the two conditional means. The exclusion
of \(r\) from its cluster's clone choices changes the allowed
placements there, but not their independence from the remaining
randomness.

\emph{A permutation-image step: the swap bijection.}
All clone placements are now fixed. Suppose the next image to
be revealed is \(\sigma_w(u)\). Let \(R\) be the set of unused
images at \(w\), and choose distinct \(z,z'\in R\). Conditional
on the previous images and on \(\sigma_w(u)=z\), the remaining
assignment is a uniform bijection onto \(R\setminus\{z\}\).
Given such a completed permutation, let
\(u'=\sigma_w^{-1}(z')\). The neighbor \(u'\) is unexposed and
different from \(u\). Swap the images at \(u,u'\), obtaining
\[
 \sigma'_w(u)=z',\qquad \sigma'_w(u')=z,\qquad
 \sigma'_w(a)=\sigma_w(a)\quad(a\notin\{u,u'\}).
\]
No previously exposed image changes, because neither \(z\) nor
\(z'\) was used in the history. The same swap is its inverse.
It is therefore a bijection from the completions with next
image \(z\) to those with next image \(z'\), and sends the
uniform law on the first set to the uniform law on the second.
Couple every other local permutation identically.

The corresponding state permutations differ only at
\((u,w)\) and \((u',w)\). Since all \(A_s,B_s\) are fixed,
at most these two summands in \(Z\) change, each by at most one.
If \(Z'\) is the value in the paired completion, it follows that
\[
 |h_i(z)-h_i(z')|
 =|\mathbb E[Z-Z'\mid\text{the fixed history and next image }z]|
 \le\mathbb E[|Z-Z'|\mid\text{the same data}]\le2.
\]
If only one image is available, the revelation is deterministic
and has conditional range zero.

\emph{Conditional ranges and the tail constant.}
At any step, let \(m_i,M_i\) be the minimum and maximum of
\(h_i(q)\) over its compatible next values. The preceding
couplings give \(M_i-m_i\le2\). Conditional on
\(\mathcal F_{i-1}\), the martingale difference
\(D_i=Y_i-Y_{i-1}\) has mean zero and lies in the interval
\[
 [m_i-Y_{i-1},\,M_i-Y_{i-1}],
\]
whose length is at most two. Thus conditional Hoeffding's
lemma, with the deterministic bounds \(c_i=2\), gives
\[
 \mathbb E[e^{\lambda D_i}\mid\mathcal F_{i-1}]
 \le e^{\lambda^2c_i^2/8},\qquad
 \sum_{i=1}^Jc_i^2=4\bigl(v+n(d-1)\bigr)\le8v.
\]
Iterating conditional expectations bounds
\(\mathbb E e^{-\lambda(Z-\mathbb EZ)}\) by
\(e^{\lambda^2v}\). Markov's inequality and the choice
\(\lambda=t/(2v)\) consequently yield, for every \(t>0\),
\[
 \Pr[Z-\mathbb EZ\le-t]\le
 \exp\!\left(-\frac{t^2}{4v}\right).
\]
Finally, \eqref{eq:14} gives \(\mathbb EZ\ge\eta v\).
The event \(Z<\eta v/2\) is therefore contained in
\(\{Z-\mathbb EZ\le-\eta v/2\}\). Taking \(t=\eta v/2\) gives

\begin{equation}\label{eq:15}
\Pr[Z<\eta v/2]
       \le\exp\left(-\frac{2(\eta v/2)^2}{8v}\right)
       =\exp(-\tau v),\qquad \tau:=1/640000.
\end{equation}
The bounds are on conditional ranges, not merely on absolute
martingale increments. This accounts for the constant in \eqref{eq:15}.
Exposing a local permutation as a single block would lose a factor
\(d\) and would not justify the conditioning argument at this degree.

The swap coupling was used only before conditioning on \(E\);
it need not preserve the single-cycle event. We now use
\(\Pr[A\mid E]\le\Pr[A]/\Pr[E]\), rather than assert the same
coupling under the conditioned law. By \eqref{eq:9},
\begin{equation}\label{eq:16}
\Pr[Z<\eta v/2\mid E]
       \le v\exp\left[-\left(\tau-\frac{27}{50d^2}\right)v\right].
\end{equation}
For every \(d\ge734\),
\begin{equation}\label{eq:17}
\frac{27}{50d^2}<\frac1{960000}
                                      =\frac{2\tau}{3}.
\end{equation}
Thus \eqref{eq:16} is at most \(v e^{-\tau v/3}\).
For \(v\ge2^{32}\), \(\log v\le v/3840000=\tau v/6\);
this follows at the endpoint from \(2^{32}>32\cdot3840000\),
and from the decrease of \(\log v/v\) thereafter. Consequently
\begin{equation}\label{eq:18}
\Pr[Z<\eta v/2\mid E]\le e^{-v/3840000}.
\end{equation}
No independence after conditioning on \(E\) is asserted.
Removing the single root stage loses at most one term. For
\(v\ge800\), \eqref{eq:18} leaves at least \(\eta v/4=v/800\)
honoured active stages, except with the displayed failure probability.

\subsection{Coefficient transfer and shared gates}\label{sub:transfer}
\subsubsection*{Shared-gate decomposition}
We give a first-principles proof of the balanced-product argument and
its coefficient charge. These are the structural and rectangular
decomposition tools used in \cite[Theorem~2 and Lemma~9]{CDGM2022};
they are not claimed as new lower-bound methods.
Write \(A\preceq B\) when every coefficient of \(A\) is at most the
corresponding coefficient of \(B\).

\begin{lemma}[Balanced products with shared gates]\label{lem:balanced}
Let \(X_1,\ldots,X_g\) be disjoint nonempty finite variable groups,
where \(g\ge4\). Suppose a monotone arithmetic circuit with \(s\)
binary arithmetic gates computes a nonzero full-group set-multilinear
polynomial \(p\). Then
\[
 p=\sum_{j=1}^t U_jV_j,\qquad t\le s,
\]
where \(U_j,V_j\) have nonnegative coefficients and are set-multilinear
on complementary group sets \(S_j,[g]\setminus S_j\), with
\(g/3\le |S_j|\le2g/3\). In particular, \(U_jV_j\preceq p\).
Constants may be arbitrary nonnegative reals, and fan-out is unrestricted.
No condition on the intermediate polynomials is assumed.
\end{lemma}

\begin{proof}
\emph{Productive gates have fixed group sets.}
First propagate zero values and discard branches that cannot contribute
to the output. A multiplication with a zero child is zero; a zero
summand can be discarded. This introduces no arithmetic gates.
Every remaining gate has a nonzero polynomial and a path to the
output whose sibling at each multiplication is nonzero. Call such a
gate productive, and write \(P_u\) for its polynomial.

Fix one such path from \(u\) to the output. At each multiplication
on the path, choose a monomial of its sibling with positive
coefficient. Addition preserves coefficientwise domination, and
multiplication by a nonnegative polynomial preserves it as well.
Induction along the path therefore gives
\[
 \lambda\,mP_u\preceq p
\]
for a monomial \(m\) and a real \(\lambda>0\). This remains true
if a sibling also depends on \(u\): its chosen monomial is a
monomial of its actual polynomial, and the same coefficientwise
inequality applies.

For every monomial \(a\) of \(P_u\), the monomial \(ma\) must thus
use exactly one variable from every group, with exponent one.
Consequently \(m\) uses a fixed subset of groups once, and every
monomial of \(P_u\) uses exactly its complementary subset once.
Denote the latter subset by \(S(u)\).
At an addition gate the nonzero children have the same group set
as the parent. At a multiplication gate the two group sets are
disjoint and their union is the parent's: multiply any nonzero
monomial of each child to obtain these assertions.

\emph{A balanced arithmetic gate exists.}
Start at the output, whose group set has size \(g\). While the
current size exceeds \(2g/3\), at an addition choose a nonzero
child; at a multiplication choose a child whose group set has
size at least \(g/3\). Such a multiplication child exists because
the two disjoint sizes sum to more than \(2g/3\).
Continue if the chosen child still has size greater than \(2g/3\).
Acyclicity and the fact that an input uses at most one group force
this descent to encounter a gate \(u\) with
\(g/3\le |S(u)|\le2g/3\). Since \(g\ge4\), this gate uses at
least two groups and is therefore an arithmetic gate.

\emph{Substitute at every use of that gate.}
Write \(U=P_u\). Replace the value computed at node \(u\) by
one fresh indeterminate \(z\), at all outgoing uses simultaneously.
Evaluate the resulting DAG formally to obtain
\[
 Q(x,z)=\sum_{k\ge0}A_k(x)z^k,\qquad A_k\in\mathbb R_{\ge0}[X].
\]
Substituting \(z=U(x)\) restores the circuit, so \(Q(x,U)=p\).
The polynomial depends on \(z\): the productive path from \(u\)
propagates a monomial containing \(z\), since all multiplication
siblings remain nonzero polynomials after the replacement.
Their substitution at \(z=U\) gives their original nonzero values.

If some \(A_k\ne0\) for \(k\ge2\), choose monomials \(a\) of
\(A_k\) and \(b\) of \(U\) with positive coefficients.
Then \(ab^k\) occurs in \(Q(x,U)=p\) with positive coefficient.
It repeats every group of the nonempty set \(S(u)\), contradicting
full-group set-multilinearity. Hence \(A_k=0\) for \(k\ge2\),
and \(A_1\ne0\). We have the exact identity
\[
 p=UA_1+A_0,\qquad A_0=Q(x,0).
\]
For any monomial \(a\) of \(A_1\) and any fixed monomial \(b\)
of \(U\), the product \(ab\) occurs in \(p\). Therefore \(A_1\)
is set-multilinear on exactly \([g]\setminus S(u)\).
Nonnegativity gives \(UA_1\preceq p\).

\emph{Remove nodes permanently and telescope.}
Apply the preceding argument to the current nonzero residual,
starting with \(R_0=p\). At iteration \(j\), take \(U_j\) to
be the selected node's value in that residual circuit, and take
\(V_j\) to be the coefficient of \(z\) after global replacement.
Permanently set the selected node to zero. The resulting circuit
computes \(R_j\), with
\[
 R_{j-1}=U_jV_j+R_j,\qquad
 0\preceq R_j\preceq R_{j-1}\preceq p.
\]
Every nonzero residual is again full-group set-multilinear, so
the argument can be repeated. Each iteration selects a different
original arithmetic node: a node permanently set to zero cannot
be productive again. Pruning does not add nodes. Thus there are
at most \(s\) iterations, and the final residual is zero; otherwise
the balanced descent would supply another eligible node.
Telescoping proves the stated decomposition and coefficient domination.

The count is of removed DAG nodes. No bound on the number of
parse trees, paths, or uses of a node is needed. All their
multiplicities are already included in the coefficient \(V_j=A_1\).
Neither \(V_j\) nor the displayed decomposition is required to
have a separate small-circuit representation.
\end{proof}

\begin{proposition}[Coefficient charge for a signed measure]\label{prop:charge}
Let \(\Omega_i=X_i\) index the choices in group \(i\), and put
\(\Omega=\prod_{i=1}^g\Omega_i\). Let \(\mu:\Omega\to\mathbb R\)
be a signed weight function such that, for every balanced partition
\(S,[g]\setminus S\) and every rectangle \(A\times B\) in that
partition,
\[
 \left|\sum_{\nu\in A\times B}\mu(\nu)\right|\le\gamma.
\]
For a full-group polynomial define
\(\mathcal M_\mu(p)=\sum_{\nu\in\Omega}\mu(\nu)[m_\nu]p\).
If \(p\) satisfies Lemma~\ref{lem:balanced} and its coefficients
are at most \(C\), then
\[
 |\mathcal M_\mu(p)|\le sC\gamma.
\]
No product-distribution assumption on \(\mu\) is made.
\end{proposition}

\begin{proof}
Fix one product \(UV\) in the lemma, with complementary group
sets \(S,T\). Let \(u_\alpha,v_\beta\) be its factors' coefficient
vectors over \(\Omega_S,\Omega_T\), including zero coefficients.
If \(UV=0\), its charge is zero, so assume both factors are nonzero.
Disjoint groups give
\[
 [m_{(\alpha,\beta)}]UV=u_\alpha v_\beta.
\]
Because the domain is the full Cartesian product and \(UV\preceq p\),
\[
 \left(\max_\alpha u_\alpha\right)
 \left(\max_\beta v_\beta\right)
 =\max_{\alpha,\beta}u_\alpha v_\beta\le C.
\]
This bounds the product of the two maxima; neither individual
maximum is assumed bounded by \(C\).

List the distinct positive levels of \(u\) as
\(0=a_0<a_1<\cdots<a_r\), and those of \(v\) as
\(0=b_0<b_1<\cdots<b_\ell\). Define
\[
 A_i=\{\alpha:u_\alpha\ge a_i\},\qquad
 B_j=\{\beta:v_\beta\ge b_j\},\qquad
 \lambda_{ij}=(a_i-a_{i-1})(b_j-b_{j-1})>0.
\]
The identities
\[
 u_\alpha=\sum_{i=1}^r(a_i-a_{i-1})
                  {\bf1}_{\{u_\alpha\ge a_i\}},\qquad
 v_\beta=\sum_{j=1}^\ell(b_j-b_{j-1})
                  {\bf1}_{\{v_\beta\ge b_j\}}
\]
give the following entrywise equality on the full Cartesian domain:
\[
 u_\alpha v_\beta
 =\sum_{i=1}^r\sum_{j=1}^\ell
       \lambda_{ij}{\bf1}_{A_i\times B_j}(\alpha,\beta).
\]
Each \(A_i\times B_j\) is a rectangle for the same balanced
partition \(S,T\). Telescoping the level increments gives exactly
\[
 \sum_{i,j}\lambda_{ij}
 =\left(\sum_i(a_i-a_{i-1})\right)
  \left(\sum_j(b_j-b_{j-1})\right)
 =a_rb_\ell=(\max u)(\max v)\le C.
\]
Thus it is the total scalar weight, regardless of how many
rectangles occur, that controls the charge:
\[
 |\mathcal M_\mu(UV)|
 =\left|\sum_{i,j}\lambda_{ij}
                  \sum_{\nu\in A_i\times B_j}\mu(\nu)\right|
 \le\gamma\sum_{i,j}\lambda_{ij}\le C\gamma.
\]
This is the rectangular decomposition of
\cite[Lemma~9, pp.~39:10--39:11]{CDGM2022}, followed by the
signed-measure estimate in its Lemma~7. The displayed proof
allows arbitrary nonnegative real coefficients and an arbitrary
signed weight function \(\mu\).
Sum this inequality over the at most \(s\) products in
Lemma~\ref{lem:balanced}.
\end{proof}

\subsubsection*{Discrepancy and coefficient transfer}

Fix a balanced partition, condition on the layout, and fix all
unhonoured bits. By locality, the pullback of any communication
rectangle is an inner-product rectangle in the remaining bits.
For the \(2^z\)-by-\(2^z\) inner-product sign matrix \(A\),
the identity \(AA^\top=2^z I\) gives norm \(2^{z/2}\).
For rectangle indicator vectors, Cauchy--Schwarz after normalization
by \(2^{2z}\) bounds the signed mass by \(2^{-z/2}\).
Using \eqref{eq:18} and \(z\ge v/800\), the original distribution has
\begin{equation}\label{eq:19}
\gamma_0\le e^{-v/3840000}+2^{-v/1600}.
\end{equation}
This is the same distribution for every balanced partition.
No union bound over partitions or simultaneously good layout
is needed.

Every layout has \(k=v-1\) bits of inner product. Its two answer
masses are \((1-2^{-k})/2\) and \((1+2^{-k})/2\).
Balance them exactly, changing total variation by \(2^{-k-1}\).
Explicitly, if \(p_i=\Pr_\Delta[g_H=i]\), set
\(\Delta_b(\nu)=\Delta(\nu)/(2p_{g_H(\nu)})\).
Then both answers have mass \(1/2\).
Every rectangle's signed expectation changes by at most \(2^{-k}\).
For \(v\ge2^{32}\), the centered maximum discrepancy is therefore
\begin{equation}\label{eq:20}
\gamma\le 2^{-cv},\qquad
                         c:=1/7680000.
\end{equation}
Indeed the left side is at most \(3\cdot2^{-v/3840000}\),
and the factor three is absorbed by half the exponent.

Define \(M(m_\nu)=\Delta_b(\nu)(2g_H(\nu)-1)\). Then
\[
             M(F_H)=0,\qquad M(P_H)=1/2,\qquad
             |M(F_H\pm\epsilon P_H)|=\epsilon/2.
\]
To make the coefficient bound in Proposition~\ref{prop:charge}
explicit here, put
\(\mu(\nu)=\Delta_b(\nu)(2g_H(\nu)-1)\) and
\(p=F_H\pm\epsilon P_H\). For each balanced product
\(U_jV_j\preceq p\) supplied by Lemma~\ref{lem:balanced},
every coefficient satisfies
\[
 0\le u_\alpha v_\beta
 \le [m_{(\alpha,\beta)}]p
 \le 1+\epsilon\le\frac32.
\]
The level-set rectangles in Proposition~\ref{prop:charge} therefore
have total weight
\[
 \sum_{i,k}\lambda_{ik}
 =(\max_\alpha u_\alpha)(\max_\beta v_\beta)
 =\max_{\alpha,\beta}u_\alpha v_\beta\le\frac32.
\]
Each rectangle has signed mass at most \(\gamma\) in absolute
value by \eqref{eq:20}, so \(|M(U_jV_j)|\le(3/2)\gamma\).
The bound holds even if one factor has very large coefficients:
only the product of the two maxima is used. Summing over
\(t\le s\) products gives
\[
 \frac{\epsilon}{2}=|M(p)|
 \le\sum_{j=1}^t|M(U_jV_j)|
 \le\frac32\,t\gamma\le\frac32\,s\gamma,
\]
and hence
\begin{equation}\label{eq:21}
s\ge\frac{\epsilon}{3\gamma}.
\end{equation}
Since \(N<Kv\), \(\epsilon\ge2^{-aN}\) with \(a=c/(2K)\)
implies \(\epsilon\ge2^{-cv/2}\). Thus
\[
                  s\ge2^{cv/2}/3\ge2^{cv/3}\ge2^{bN}
\]
with \(b=c/(3K)\); the middle inequality is already valid under
the stated cutoff. This proves Theorem~\ref{thm:finite}.


\section{Explicit families and actual variable counts}\label{sec:explicit}

\subsection{An explicit LPS family}\label{sub:lps}
Fix the prime \(p=733\). For every field prime
\[
                     q\equiv1\pmod{2932},\qquad q>1466,
\]
the standard LPS graph \cite{LPS1988} \(B_q\) is the nonbipartite Cayley graph
on \(\operatorname{PSL}_2(\mathbb F_q)\) of degree \(d=p+1=734\).
It has
\[
 n=\frac{q(q^2-1)}2,\qquad v=367q(q^2-1).
\]
Indeed \(p,q\equiv1\pmod4\), and quadratic reciprocity plus
\(q\equiv1\pmod p\) gives \((p/q)=1\).
The LPS spectral theorem gives
\[
 \max_{i\ge2}|\theta_i|\le\frac{2\sqrt{733}}{734}<\frac2{27}.
\]
Connectivity and nonbipartiteness hold on this PSL component.
The conservative bound \(q>2p\) ensures the standard generators
remain nonidentity and distinct modulo \(q\), so the graph is simple.
No girth theorem is used.

One can construct the fixed generator list from integer
quaternions \((a_0,a_1,a_2,a_3)\) of norm \(733\), where \(a_0\)
is positive and odd and the other coordinates are even. With
\(i^2=-1\) modulo \(q\), send it to the projective matrix
\[
 \begin{pmatrix}
   a_0+i a_1&a_2+i a_3\\
   -a_2+i a_3&a_0-i a_1
 \end{pmatrix}.
\]
Its determinant is \(733\), a square modulo \(q\); normalizing
to determinant one gives a PSL element independent of the sign
of the square root. This is the conventional LPS construction.
For the fixed \(p\), exhaustive enumeration of the generator
list and of determinant-one matrices modulo \(\{\pm I\}\)
takes time polynomial in the graph size.

For clarity, simplicity can also be checked directly in this
single-prime construction. Normalize every matrix by the same
square root of \(733\). Equality in PSL then means equality
or negated equality of the two normalized matrices. The quaternion
matrix map is invertible over \(\mathbb F_q\), so the coordinate
vectors are congruent up to the same sign. Since every coordinate
has absolute value at most \(\sqrt{733}\) and \(q>2\cdot733\),
the congruence is an integer equality. The positive \(a_0\)
convention eliminates the minus case. An identity generator
would require \(a_1=a_2=a_3=0\), contradicting \(a_0^2=733\).
Quaternion conjugation gives inverse closure of the generator list.

Take \(L=32\). The exact host degree and variable factor are
\[
 D_H=17263679,\qquad K=LD_H=552437728.
\]
With \(c=1/7680000\), the constants in the theorem can be
chosen as the positive rationals
\[
 a=\frac1{8485443502080000},\qquad
 b=\frac1{12728165253120000}.
\]
The cutoff \(v\ge2^{32}\) is already implied by \(q>1466\).
All these constants are deliberately conservative.


\subsection{Every sufficiently large variable count}\label{sub:padding}
The core graph theorem gives an explicit infinite subsequence.
We give a padding argument for all sufficiently large
actual variable counts. It uses only the prime number theorem in
the fixed progression \(1\pmod{2932}\), in addition to the core
graph result.

Set \(m=D_H\), \(A=367L\). The core number of groups is
\[
                    g_q=Aq(q^2-1)-1.
\]
For any sufficiently large integer \(N\), let
\(g=\lfloor N/m\rfloor\) and \(x=(g/A)^{1/3}\).
The prime number theorem in this fixed progression gives
an admissible prime \(q\in(x/2,x]\) once \(x\) is large enough;
for a stronger quantitative statement see \cite{ThornerZaman2021}. Then \(g_q\le g\). Since \(q\ge2\) and \(q^2-1\ge3q^2/4\),
\[
            g_q\ge\frac{3g}{32}-1\ge\frac g{16}
                     \qquad(g\ge32).
\]
The core variable count \(N_0=mg_q\) is therefore at least
\(N/32\), using \(mg\ge N/2\) for \(N\ge2m\).

Append \(g-g_q\) fresh groups of size \(m\), and, if \(N-mg>0\),
one final fresh group of size \(N-mg\). Let \(G\) be the product
of all appended group sums, with \(G=1\) if none are needed.
Define
\[
             F_N=F_{H_q}G,\qquad P_N=P_{H_q}G.
\]
These are explicit polynomials on exactly \(N\) variables.
Fixing one variable in each appended group to one and the
others to zero restricts \(F_N\pm\epsilon P_N\) to the core
target without increasing monotone circuit size.

For \(a'=a/32\) and \(b'=b/32\), the conditions
\[
                       2^{-a'N}\le\epsilon\le1/2
\]
imply \(\epsilon\ge2^{-aN_0}\), since \(N_0\ge N/32\).
The core theorem therefore gives size at least
\(2^{bN_0}\ge2^{b'N}\).
The finitely many smaller indices may be filled arbitrarily.
Primality testing or direct enumeration in this fixed progression
is polynomial in \(N\); the determinant construction stays uniform.

This does not assert a universal discrepancy estimate for every
balanced partition of the padded variables. The lower bound is
transferred by an arithmetic restriction after the shared-gate
theorem has already been proved on the core.

This completes the proof of Theorem~\ref{thm:asymptotic},
with its constants taken to be \(a',b'\).


\section{Full support and sharp approximation}\label{sec:approx}

\subsection{Two coefficient values}\label{sub:twolevels}
The polynomial \(R_H=2F_H-P_H\) has every full-group monomial,
each with coefficient one or two. It is in uniform VP over the
integers. A monotone circuit of size \(s\) for \(R_H\), followed
by one multiplication by \(1/2\), computes
\(F_H-\tfrac12P_H\). Thus \(s+1\ge2^{bN}\), implying
\(s=2^{\Omega(N)}\). The padded family has the same conclusion.

More generally, let \(Q=F_H\pm\epsilon P_H\). Suppose
\(\widetilde Q\) is a nonnegative set-multilinear polynomial
on exactly the full set of tail groups and
\[
 \big|[m_\nu]\widetilde Q-[m_\nu]Q\big|\le\epsilon/4
                         \quad\text{for every neighbor map }\nu.
\]
The centered coefficient functional has total absolute weight
one, so \(|M(\widetilde Q)|\ge\epsilon/4\).
Its coefficients are at most \(1+\epsilon+\epsilon/4<2\).
The same balanced-product argument gives
\[
            \operatorname{monSIZE}(\widetilde Q)
                      \ge\frac{\epsilon}{8\gamma}.
\]
This is a strongly exponential coefficient-approximation lower
bound in the same perturbation range, after decreasing the
exponent constant if needed. The requirement that the output
be set-multilinear on exactly these groups excludes uncontrolled
extra monomials; it is not an approximation claim for arbitrary
nonhomogeneous output.

Conversely, the number of full-group monomials is
\(D_H^{Lv-1}=2^{O(N)}\). Listing them gives a monotone upper
bound \(2^{O(N)}\) for the exact targets. The size scale is
therefore \(2^{\Theta(N)}\) with fixed, different constants.
\subsection{Exact coefficient-approximation tradeoff}

Use the same host, full product \(F_H\), arborescence polynomial
\(P_H\), and exactly centered distribution \(\Delta_b\).
Let \(\gamma\) bound its rectangle discrepancy for every balanced
partition. The Eulerian theorem gives \(\gamma\le2^{-cv}\).
The signed coefficient functional
\[
 M(Q)=\sum_\nu\Delta_b(\nu)(2g_H(\nu)-1)[m_\nu]Q
\]
assigns total positive weight \(1/2\) and total negative weight
\(-1/2\).

Fix \(0<\ell<u\). Consider either assignment of the two
coefficient levels to the two answers:
\[
 R_{\ell,u}=\ell P_H+u(F_H-P_H)
 \quad\text{or}\quad
 R_{\ell,u}=uP_H+\ell(F_H-P_H).
\]
Let \(1\le\kappa<\sqrt{u/\ell}\). Suppose a monotone circuit
computes a full-group set-multilinear polynomial \(\widetilde R\)
whose coefficients satisfy
\begin{equation}\label{eq:A1}
\frac{[m_\nu]R_{\ell,u}}{\kappa}
       \le [m_\nu]\widetilde R
       \le \kappa[m_\nu]R_{\ell,u}
                          \qquad\text{for every }\nu.
\end{equation}
Then its size satisfies the explicit bound
\begin{equation}\label{eq:A2}
\boxed{\quad
    s\ge\frac{u-\kappa^2\ell}{2\kappa^2u\,\gamma}.
       \quad}
\end{equation}

\emph{Proof.} First take the assignment with level \(\ell\) on
\(g_H=1\) and level \(u\) on \(g_H=0\). Exact centering and
\eqref{eq:A1} give
\[
 \begin{aligned}
 M(\widetilde R)
 &\le \frac{\kappa\ell}{2}-\frac{u}{2\kappa}\\
 &=-\frac{u-\kappa^2\ell}{2\kappa}<0.
 \end{aligned}
\]
The other assignment reverses the sign and gives the same
lower bound on \(|M(\widetilde R)|\).
Every output coefficient is at most \(\kappa u\).
The shared-gate balanced-product decomposition therefore writes
the output as at most \(s\) terms, each having rectangle-mixture
mass at most \(\kappa u\). Each term's functional magnitude is
at most \(\kappa u\gamma\). Comparing the two bounds yields
\eqref{eq:A2}. No precision restriction on the constants is used. \(\square\)

For fixed \(\ell,u,\kappa\) with \(\kappa<\sqrt{u/\ell}\), the
factor preceding \(1/\gamma\) is a positive constant.
Consequently \eqref{eq:A2} is \(2^{\Omega(N)}\) in actual variables.
If \(\kappa\) depends on \(N\) and approaches the threshold,
only the explicit tradeoff \eqref{eq:A2}, rather than a uniform exponent,
is being asserted.

\subsection{Matching easy approximation at the threshold}

At \(\kappa_0=\sqrt{u/\ell}\), the easy polynomial
\begin{equation}\label{eq:A3}
\sqrt{\ell u}\,F_H
\end{equation}
satisfies \eqref{eq:A1} for both coefficient levels. For the lower level,
\(\sqrt{\ell u}=\kappa_0\ell\); for the higher level,
\(\sqrt{\ell u}=u/\kappa_0\).
Thus \eqref{eq:A3} has an \(O(N)\)-size monotone circuit and reaches
the threshold exactly in the real-constant model.

In particular, for
\[
                              R_H=2F_H-P_H,
\]
the symmetric multiplicative coefficient-approximation threshold
is \(\sqrt2\): for every fixed \(1\le\kappa<\sqrt2\),
\begin{equation}\label{eq:A4}
s\ge\frac{2-\kappa^2}{4\kappa^2\gamma}
                         =2^{\Omega(N)},
\end{equation}
whereas \(\sqrt2 F_H\) is an \(O(N)\)-size \(\sqrt2\)-approximation.
The lower bound is for unrestricted monotone circuits with
shared gates, not just sums of products.

\subsection{Uniform pointwise approximation needs no group-structure assumption}

The same sharp threshold holds if an arbitrary monotone circuit
is required to approximate \(R_{\ell,u}\) pointwise even on the
bounded cube:
\begin{equation}\label{eq:A5}
\frac{R_{\ell,u}(x)}{\kappa}\le \widetilde R(x)
                   \le\kappa R_{\ell,u}(x)
                       \qquad\text{for every }x\in[0,1]^N.
\end{equation}
No degree or set-multilinearity assumption on the circuit's output
or intermediate gates is needed. The following standard extraction argument
\cite{JerrumSnir1982,Jukna2016} preserves size even when gates are shared.

\begin{lemma}[Least homogeneous component]\label{lem:least} If a nonzero polynomial \(Q\) is
computed by a size-\(s\) monotone circuit, its nonzero homogeneous
component of least total degree has a monotone circuit of size
at most \(s\).
\end{lemma}

\begin{proof} Prune zero and unused gates. Assign each remaining gate
its least monomial degree. At addition this is the minimum of
the children's degrees; at multiplication it is their sum,
because no cancellation is possible.
Build one corresponding gate for each original gate. At an
addition with unequal least degrees retain only the child of
smaller degree; at equal degrees retain their addition.
At multiplication retain the product of the two least components.
Inputs remain variables or nonnegative constants.
Induction computes each gate's least homogeneous component.
Each original gate contributes at most one arithmetic gate,
and all of its uses share that same constructed gate.
Thus size does not increase.\end{proof}

Let \(g\) be the number of tail groups, so the target
\(R_{\ell,u}\) is homogeneous of degree \(g\) and positive
at \(\mathbf1\). Let \(d\) be the least degree of
\(\widetilde R\). Substitute \(x=t\mathbf1\) in \eqref{eq:A5},
divide by \(t^g\), and let \(t\) decrease to zero.
If \(d<g\), the upper bound fails; if \(d>g\), the lower
bound fails. Hence \(d=g\).

Let \(Q_g\) be that least homogeneous component. It has
monotone size at most the size of the original approximator.
For every fixed \(x\in[0,1]^N\), the points \(tx\) remain
in the cube for \(0<t\le1\). Dividing \eqref{eq:A5} at \(tx\) by
\(t^g\) and taking the limit gives
\[
          R_{\ell,u}(x)/\kappa\le Q_g(x)\le\kappa R_{\ell,u}(x).
\]
On a face where every variable of one tail group is zero,
the target vanishes. Nonnegativity therefore forces every
monomial of \(Q_g\) to meet every tail group. Its degree is
exactly \(g\), so it uses each group exactly once.
Thus \(Q_g\) is full-group set-multilinear.

Evaluate at the one-hot point for each neighbor map: set
its selected variable in every group to one and all others
to zero. This isolates the corresponding coefficient of
\(Q_g\), proving \eqref{eq:A1} for \(Q_g\). Applying \eqref{eq:A2} to its
no-larger circuit proves the same bound for the original
arbitrary monotone approximator.

Conversely, a full-group polynomial satisfying \eqref{eq:A1} satisfies
the pointwise inequalities on the whole nonnegative orthant
by summing termwise. In particular \eqref{eq:A3} attains the threshold
on the cube and on the orthant. The optimal monotone size
for coefficient approximation equals that for uniform cube
approximation, with no size loss in either direction.
The latter assertion compares the optimal circuit sizes;
it does not say that every cube approximator itself is
set-multilinear.

For example \(xy+(xy)^2\) is a factor-two approximation to
\(xy\) on \([0,1]^2\) and is not multilinear. The least-component
operation extracts \(xy\). This explains why forcing the
original output to be multilinear would have been too strong,
while the size lower bound still transfers.

The same conclusion holds if approximation is required on the
open positive cube, since polynomial continuity supplies its
boundary. It also holds on the entire nonnegative orthant.
The argument uses the continuum of scales tending to zero;
it does not give an approximation theorem on Boolean inputs alone.

\subsection{Padding and scope}

The all-size padding in Section~\ref{sub:padding} preserves this
corollary. Choosing one variable in each added group to be one
and the others zero isolates one padded extension of each core
monomial. Coefficient approximation restricts to coefficient
approximation on the core; pointwise approximation restricts
directly. A constant fraction of the actual variables remains.

The two-level target has a polynomial-size general arithmetic
circuit because \(P_H\) is a rooted determinant. At levels one
and two, no growing coefficient precision is needed even in
the exact target. This is a sensitive monotone arithmetic
separation and its sharp approximation corollary.

The Boolean support function of \(R_H\) is the same easy support
function as \(F_H\). This result does not imply a monotone Boolean
circuit lower bound through a support conversion. It also does
not supply Hrube\v{s}'s additional hypotheses for general arithmetic
lower bounds. 


\section{Exact counting on the Boolean cube}\label{sec:counting}

Exact counting requires equality with the target on every Boolean
assignment, in the sense of \cite{Jukna2016}. The full-group normalization
below is related to the consistent-to-decomposable normalization
of sum-product networks \cite{Peharz2015,MartensMedabalimi2015}.
We include its proof to specify the hypotheses, arbitrary-real-constant
model, and cost in shared arithmetic gates.

\subsection{Counting normalization statement}

Partition the variables into \(g\ge1\) nonempty groups
\(X_1,\ldots,X_g\). Let \(p\ne0\) be a polynomial with nonnegative
coefficients whose every monomial uses exactly one variable from
each group, with exponent one. Suppose a size-\(s\) monotone
arithmetic circuit computes some polynomial \(Q\) satisfying
\begin{equation}\label{eq:CN1}
Q(x)=p(x)\quad(x\in\{0,1\}^N).
\end{equation}
The circuit has binary \(+\) and \(\times\) gates, arbitrary
nonnegative real constants, and unrestricted reuse; \(Q\) need
not be multilinear or homogeneous.

\begin{theorem}[Full-group counting normalization]\label{thm:counting} There is a monotone circuit computing
\(p\) as a formal polynomial, of size at most
\begin{equation}\label{eq:CN2}
(2g+1)s+g.
\end{equation}
The transformation preserves shared subcircuits and takes
polynomial time in the original circuit and variable count.
It does not depend polynomially on the potentially huge formal
degree of \(Q\).
\end{theorem}

\subsection{Profiles of productive gates}

For a polynomial \(U\), write \(\operatorname{lin}(U)\) for
the result of replacing each positive exponent by one and
combining like terms. A multilinear polynomial is uniquely
determined by its values on the Boolean cube; this follows
inductively by writing it as \(A+x_NB\) and evaluating at
\(x_N=0,1\). Thus \eqref{eq:CN1} implies
\begin{equation}\label{eq:CN3}
\operatorname{lin}(Q)=p.
\end{equation}
Nonnegativity then implies that every monomial of \(Q\) has
exactly one distinct variable from each group, although that
variable may have a higher exponent.

Prune zero and unused gates. For each productive gate \(v\),
let \(S_i(v)\subseteq X_i\) be the set of variables from group
\(i\) appearing in any monomial at \(v\). Call \(i\) flexible
at \(v\) if \(|S_i(v)|\ge2\), and fixed at \(v\) if
\(|S_i(v)|=1\). Empty profiles are allowed.

Every monomial at \(v\) is a partial choice: it cannot contain
two distinct variables from a group. Indeed choose a nonzero
output parse context containing an occurrence of \(v\), and
fix monomials in all its other factors and occurrences.
Every monomial at the designated occurrence, times this
fixed context, occurs at the output with a positive coefficient.
Two different variables of a group could not disappear.

Moreover, every flexible group at \(v\) occurs in every
monomial at \(v\). For the same fixed context, if it contained
a variable from such a group, some monomial at \(v\) with a
different variable would create an invalid output support.
Hence the context contains no variable from that group.
The full-group output condition then forces every monomial
at \(v\) to supply it.

At a multiplication gate \(v=u t\), if both children have a
nonempty profile in a group, both profiles must be the same
singleton. Otherwise a cross-product of monomials would
contain two different variables from that group. Consequently
the flexible groups at \(v\) are the disjoint union of the
flexible groups at its two children.

At an addition gate \(v=u+t\), profiles are their unions.
If a group is flexible at \(v\) but not flexible at a
nonzero child, that child has a singleton profile there,
and its variable occurs in every child monomial. An empty
profile or an omitting monomial would contradict the
mandatory presence of every flexible group at \(v\).

\subsection{Constructing the normalized circuit}

Define \(J_v\) by setting all variables of fixed groups at
\(v\) to one, then linearizing the remaining polynomial.
The preceding facts show that \(J_v\) is set-multilinear on
exactly the flexible groups at \(v\).

For an input variable, \(J_v=1\). For a constant input,
\(J_v\) is that constant.

At a multiplication gate the flexible group sets are disjoint,
and every other present group is fixed at both relevant
children and parent. Hence
\begin{equation}\label{eq:CN4}
J_v=J_uJ_t.
\end{equation}
This uses one multiplication gate.

At an addition gate, let \(M_u\) be the product of the unique
child variables from groups that are flexible at \(v\) and
fixed at \(u\); define \(M_t\) analogously. Each is a product
of at most \(g\) distinct variables. The mandatory-presence
property shows that restoring one copy of each such variable
exactly reverses its removal under linearization. Therefore
\begin{equation}\label{eq:CN5}
J_v=M_uJ_u+M_tJ_t.
\end{equation}
Each restored product can be formed by multiplying its
variables into the already computed child output. This costs
at most \(2g\) multiplications and one addition per original
addition gate.

Construct one normalized output for each original gate and
reuse it at every parent. Finally, at the output, restore
one copy of the unique variable from each globally fixed
group. All these groups are mandatory by \eqref{eq:CN3}. The result
is exactly \(p\), at cost at most \(g\) more multiplications.
The total size is bounded by \eqref{eq:CN2}.

Profiles are computed from the input variables by set union
at both addition and multiplication. The latter union rule
uses the fact that all productive polynomials are nonzero
and no cancellation occurs. All other operations of the
construction are explicit. This proves the theorem.
\(\square\)

\subsection{Consequence for the sensitive VP family}

Let \(C_{01}(T)\) denote monotone circuit size required only
to agree with \(T\) on the entire Boolean cube. Apply the
normalization theorem to the full-group target
\(T=F_H\pm\epsilon P_H\). The coefficient discrepancy bound
in the Eulerian proof gives
\begin{equation}\label{eq:CN6}
C_{01}(T)\ \ge\
 \frac{\epsilon/(3\gamma)-g}{2g+1},\qquad
 g=|V(H)|-1,\qquad \gamma\le2^{-cv}.
\end{equation}
In particular, for the same fixed host parameters and
\(2^{-aN}\le\epsilon\le1/2\),
\begin{equation}\label{eq:CN7}
C_{01}(T)=2^{\Omega(N)}
\end{equation}
for sufficiently large \(N\). The polynomial normalization
factor is absorbed into the exponent. This statement does
not claim the identical numerical exponent at the original
finite cutoff for arbitrary fixed clone counts.

At \(\epsilon=1/2\), the two-coefficient target \(2F_H-P_H\)
has the same strongly exponential exact counting lower
bound. Its values on Boolean inputs count weighted full-group
monomials. This is stronger than requiring the competing
circuit to compute the formal polynomial, while still being
an arithmetic counting model, not a Boolean decision model.

The normalizer uses \eqref{eq:CN3}, which follows from exact equality
on all Boolean assignments. Approximation, or equality only
on one-hot assignments, does not supply \eqref{eq:CN3}. In particular,
an approximator may contain monomials using two distinct
variables from one group. No Boolean approximation theorem
is inferred here.


\section{Scope and further questions}\label{sec:scope}

The host polynomial has full support within its full-group monomial
basis. Replacing arithmetic addition and multiplication by Boolean OR
and AND therefore produces the same easy support predicate as for
\(F_H\). The lower bound concerns arithmetic coefficients and does
not give a Boolean circuit lower bound by a support conversion.

The distinction between complete input domains and promises is also
substantive. On the promise of exactly one selected variable in every
outgoing group, the targets have polynomial-size monotone arithmetic
circuits. To see this, represent the complement of \(x_{z,w}\) on
the promise by the sum of the other variables in its group.
Prefix and suffix sums produce all these complements in \(O(N)\) gates.
A Boolean reachability circuit can now be simulated with two
nonnegative arithmetic outputs \(a,\bar a\) for each Boolean gate.
The identities
\[
\begin{array}{lll}
\text{AND:}&ab,&\bar a+a\bar b,\\
\text{OR:}&a+\bar a b,&\bar a\bar b
\end{array}
\]
compute a gate and its complement whenever the incoming pairs are
complementary bits; NOT merely exchanges the outputs.
Reachability for every vertex uses polynomially many Boolean gates.
If \(A,\bar A\) are the resulting arborescence-indicator pair, then
\(1+\bar A\), \(1+\epsilon A\), and
\((1-\epsilon)+\epsilon\bar A\) compute respectively
\(2F_H-P_H\), \(F_H+\epsilon P_H\), and \(F_H-\epsilon P_H\)
on the promise, since \(F_H=1\) there.
These circuits need not compute the formal polynomials or their values
on the entire Boolean cube. Multiplying by \(F_H\) can enforce vanishing
on all group-zero faces without changing the promised values.

Division is another distinct model. Subtraction-free rational algorithms
for arborescence polynomials are known \cite{FGK2013}. Our lower bound
is for division-free circuits and does not preclude such algorithms.
Similarly, the perturbation interval in Theorem~\ref{thm:finite} does
not supply the additional hypotheses in Hrube\v{s}'s criteria for
general arithmetic lower bounds \cite{Hrubes2020}. In particular,
we do not obtain a separation of VP from VNP or of P from NP.

Several questions remain. Can the required base degree and clone count
be substantially reduced? Can the sensitive size bound be extended to
smaller perturbations? Does a corresponding approximation bound hold
when the domain is only the Boolean cube? Finally, can strongly
exponential sensitive hardness be achieved for a family with
polynomial-size general arithmetic formulas? The determinant upper
bound here supplies polynomial-size circuits, not that stronger
formula upper bound.

\section*{Acknowledgment}
\addcontentsline{toc}{section}{Acknowledgment}
The author acknowledges OpenAI GPT-6 for assistance with research
exploration, proof development, verification code, and manuscript
preparation. The finite computations are supplementary checks;
the asymptotic claims rest on the written arguments.



\appendix


\section{Direct BEST counting and sampling}\label{app:best}
The following standard last-exit-tree proof \cite{FarrellLevine2015} is included to
fix the counting convention. Replace every base edge by its two
directed arcs. Fix a base start vertex \(o\).
Choose a directed spanning tree toward \(o\). At each \(u\ne o\),
order its outgoing arcs with the tree arc last, and order the
other \(d-1\) arcs arbitrarily. At \(o\), choose any order of its
\(d\) outgoing arcs. Starting at \(o\), always take the next unused
outgoing arc in that vertex's order.

The walk cannot stop away from \(o\), by balance of indegree and
outdegree. If it stops at \(o\) with an unused outgoing arc
somewhere, follow unused last-exit tree arcs toward \(o\).
The first fully exhausted vertex on that tree path would have an
unused incoming arc, contradicting balance for the closed walk.
Thus all arcs are used, and the walk is an Eulerian tour.

Conversely, the last outgoing arcs of nonroot vertices in an
Eulerian tour form a tree toward \(o\). A directed cycle of these
last exits would force their exit times to increase strictly
around the cycle, which is impossible. The outgoing orders and
last-exit tree are uniquely recovered from the tour.

There are therefore
\[
                     t_o(B)\,d!\,((d-1)!)^{n-1}
\]
linear Eulerian tours starting at \(o\). Each cyclic transition
configuration is represented by exactly \(d\) such tours, one
for each possible first outgoing arc at \(o\). Dividing by \(d\)
gives \(t_o(B)((d-1)!)^n\), proving \eqref{eq:5}.

This also gives an efficient uniform sampler: choose a uniform
spanning tree of \(B\), orient it toward \(o\), choose all the
allowed outgoing orders uniformly, run the tour, and forget its
initial position. Every single-cycle transition configuration
has the same number \(d\) of preimages.
A uniform spanning tree can, for example, be sampled in expected
polynomial time by successive
edge inclusion decisions using exact matrix-tree counts and
deletion/contraction. Those counts have polynomial bit length
and require polynomially many determinant computations.
Sampling the exact rational branch probabilities takes expected
polynomially many random bits. No rejection sampling with
exponentially small acceptance is needed.


\section{A smaller physical host}\label{app:thin}
Let \(S\) be the oriented base edges and define a simple
undirected graph \(T_B\) on \(S\) by
\[
 \{(u,w),(x,z)\}\in E(T_B)
                 \quad\Longleftrightarrow\quad w=x\ \text{or}\ z=u.
\]
There are no loops because \(B\) is simple.
For \(s=(u,w)\), there are \(d\) successor states \((w,z)\)
and \(d\) predecessor states \((z,u)\). Their intersection
is exactly \(\{(w,u)\}\), so \(T_B\) has degree \(2d-1\).
It is connected, for example because any Euler tour of the
bidirected connected base traces a spanning cycle in \(T_B\).

The host
\[
                         H^\circ=T_B[K_L]
\]
has \(Lv\) physical vertices and constant degree
\[
                         D^\circ=2dL-1.
\]
Every intercluster arc in the direct gadget goes from an oriented
edge state to one of its successor states, and every other arc
stays in a state cluster. Consequently every possible gadget
assignment lies in \(H^\circ\).

The entire proof of Sections~\ref{sec:distribution} and \ref{sec:discrepancy} therefore applies
verbatim to this smaller host: its probability space, fractional
partition calculation, exposure concentration, parity identity,
total-domain rectangle argument, and shared-gate coefficient
transfer do not use any additional host edges. In particular,
with actual variable count
\[
                 N^\circ=(Lv-1)D^\circ,
\]
one may take the same \(c=1/7680000\) and
\[
 a^\circ=\frac{c}{2LD^\circ},\qquad
 b^\circ=\frac{c}{3LD^\circ}.
\]
For \(v\ge2^{32}\) and
\(2^{-a^\circ N^\circ}\le\epsilon\le1/2\),
\[
 \operatorname{monSIZE}
       (F_{H^\circ}\pm\epsilon P_{H^\circ})\ge2^{b^\circ N^\circ}.
\]

This is a direct adaptation of the proved discrepancy argument,
not a claim that deleting edges from an arbitrary hard polynomial
preserves hardness. The same hard distribution is supported in
the smaller total Cartesian domain, and that fact is what allows
the transfer to be repeated.

For \(d=734,L=32\), the degree drops from \(17263679\)
to \(46975\), while the number of physical vertices is unchanged.
The main statement keeps the simpler notation \(B[K_{dL}]\);
this optional host gives better fixed constants. Neither change
alters the asymptotic \(2^{\Theta(N)}\) scale.



\section{General conditioning and fixed-gap expanders}\label{app:general}
\subsection{The finite conditioning criterion}

Let \(B\) be a connected simple \(d\)-regular graph with \(n\)
vertices and normalized second eigenvalue at most
\(0\le\rho<1\). Put \(v=dn\), fix \(L\ge5\), and use
the direct gadget on \(H=B[K_{dL}]\). Write
\[
 \delta=\frac{1-\rho}{2},\qquad
 q_v=\frac{Lv-1}{3Lv},\qquad
 w_v=\left(\delta q_v(1-q_v)-\frac1L\right)_+^2-\frac2v,
 \qquad
 \beta=\frac1{2(1-\rho)d^2}.
\]

For every real \(t\) satisfying \(0<t<w_v\), the exactly
centered distribution of the Eulerian construction obeys
\begin{equation}\label{eq:G1}
\boxed{\quad
 \max_{\text{balanced }A}
    \operatorname{disc}_{\Delta_b,A}(g_H)
 \le
 v\exp\!\left[-\left(\frac{(w_v-t)^2}{4}-\beta\right)v\right]
       +2^{-(tv-1)/2}+2^{1-v}.
 \quad}
\end{equation}
The maximum is over partitions of whole outgoing-choice groups
with each side between one third and two thirds of the groups.
The predicate and rectangles are on the full Cartesian product
of all neighbor choices.

\emph{Proof.} Fix any such partition and assign the root either
analysis color. Its mean Alice density lies in
\([q_v,1-q_v]\). The state-transition gap is at least
\(\delta\), with \(0<\delta\le1/2\). The pair estimate in Section~\ref{sub:mean}
therefore gives
\[
                         \mathbb E Z\ge w_v v.
\]
The subtraction of two is solely the expectation loss at
the root cluster's incoming and outgoing transition.
Clones remain independent across clusters before conditioning.

Expose the \(v\) clone choices, then the local permutation
images individually. The sum of squared conditional range
lengths is at most \(8v\). Hence
\[
 \Pr[Z<tv]\le
       \exp\!\left[-\frac{(w_v-t)^2}{4}v\right].
\]
This use of the mean lower bound is valid because \(t<w_v\).
The BEST count and logarithmic trace estimate give
\[
                   \Pr[E]\ge v^{-1}e^{-\beta v}.
\]
Dividing the preceding tail by this probability proves
the first term of \eqref{eq:G1}. On a remaining layout there are
at least \(tv-1\) honoured active stages, after removing
the root stage. Fixing all other bits pulls every rectangle
back to an inner-product rectangle, proving the second term.
There are exactly \(v-1\) inner-product stages on every
layout. Exact answer centering changes signed rectangle
expectation by at most \(2^{1-v}\), giving the final term.

All these estimates hold for each fixed partition under
the same distribution. No union bound over partitions,
post-conditioning independence, or promised input domain
has been introduced. \(\square\)

As \(v\) tends to infinity with fixed parameters,
\[
            w_v\longrightarrow
            w_\infty=
             \left(\frac{1-\rho}{9}-\frac1L\right)_+^2.
\]
Thus a sufficient asymptotic condition is
\begin{equation}\label{eq:G2}
\boxed{\qquad
       \frac1{2(1-\rho)d^2}
         <\frac14
            \left(\frac{1-\rho}{9}-\frac1L\right)_+^4 .
       \qquad}
\end{equation}
Indeed choose a fixed \(t\) with
\(0<t<w_\infty-2\sqrt\beta\). Then the first exponent
in \eqref{eq:G1} is bounded away from zero for all sufficiently
large \(v\), as are the other exponential rates. Therefore
the centered discrepancy is \(2^{-\Omega(v)}\).
This is a sufficient criterion, not a converse or an optimal
degree threshold.

\subsection{Any fixed expander gap can meet the criterion}

It is useful to give a deliberately loose closed-form choice.
Write \(\kappa=1-\rho\), where \(0<\kappa\le1\). It suffices
to choose fixed integers satisfying
\begin{equation}\label{eq:G3}
L\ge\frac{18}{\kappa},\qquad
              d\ge2592\,\kappa^{-5/2}.
\end{equation}
Then \(w_\infty\ge\kappa^2/324\), so eventually
\(\mathbb E Z\ge\eta v\), where
\[
       \eta=\frac{\kappa^2}{648},\qquad
       \tau=\frac{\eta^2}{16}
            =\frac{\kappa^4}{6718464}.
\]
The degree choice gives
\[
                  \beta=\frac1{2\kappa d^2}\le\frac{\tau}{2}.
\]
Using the threshold \(\eta v/2\) in the concentration argument
therefore gives
\[
 \Pr[Z<\eta v/2\mid E]\le v e^{-\tau v/2}
                                  \le e^{-\tau v/4}
\]
for sufficiently large \(v\). The same construction has at
least \(\eta v/4\) honoured active stages outside this mass.
Consequently its centered discrepancy is at most
\begin{equation}\label{eq:G4}
e^{-\tau v/4}
                     +2^{-\eta v/8}+2^{1-v}
                     =2^{-\Omega(v)}.
\end{equation}
The implicit constants depend only on the fixed gap and
chosen degree and clone count.

\subsection{A fixed clique blow-up makes an arbitrary expander eligible}

\begin{theorem}[Fixed-gap expander families]\label{thm:general} Fix an integer \(\Delta\ge2\) and a real
\(0<\kappa_0\le1\). There is an integer \(M\), depending only
on these fixed parameters, with the following property.
For every sufficiently large connected simple
\(\Delta\)-regular graph \(G\) satisfying
\[
                      \lambda_2(A_G/\Delta)\le1-\kappa_0,
\]
the host \(H=G[K_M]\) has fixed degree and admits constants
\(a,b>0\), independent of \(|V(G)|\), such that
\begin{equation}\label{eq:G5}
\operatorname{monSIZE}(F_H\pm\epsilon P_H)\ge2^{bN}
                \qquad(2^{-aN}\le\epsilon\le1/2),
\end{equation}
where \(N=(|V(H)|-1)\deg(H)\) is its actual directed-arc
variable count. The circuit model, total communication
predicate, and coefficient definitions are exactly those
of the Theorem~\ref{thm:finite}. If \(G\) is a uniform explicit
family, the determinant upper bound is uniform as well,
with the same convention on the epsilon parameter.
\end{theorem}

\begin{proof} Put
\[
                  \kappa=\frac{\Delta\kappa_0}{\Delta+1}.
\]
Choose a fixed integer \(s\) so large that
\[
        d=s(\Delta+1)-1\ge2592\,\kappa^{-5/2},
\]
and put \(B=G[K_s]\). This is connected, simple, and
\(d\)-regular. Its adjacency matrix is
\[
                  A_B=(A_G+I)\otimes J_s-I.
\]
The subspace constant within each \(s\)-clique has normalized
eigenvalues
\[
                 \frac{s(\Delta\theta+1)-1}{d},
\]
where \(\theta\) runs over the normalized eigenvalues of \(G\).
The orthogonal subspace has normalized eigenvalue \(-1/d\).
For every nonconstant lifted eigenvector,
\[
 \frac{s(\Delta\theta+1)-1}{d}
       \le1-\frac{s\Delta\kappa_0}{d}
       \le1-\kappa.
\]
Also \(-1/d\le1-\kappa\), since \(0<\kappa<1\).
Thus \(B\) satisfies the second-eigenvalue hypothesis with
\(\rho=1-\kappa\).

Choose \(L\ge18/\kappa\). The preceding estimates provide
centered discrepancy \(2^{-cv}\) for some fixed \(c>0\),
where \(v=d|V(B)|=ds|V(G)|\), on
\[
                  H=B[K_{dL}]\cong G[K_{sdL}].
\]
Set \(M=sdL\). The blow-up factor is a fixed integer.
With \(K=L((d+1)dL-1)\), the actual count satisfies \(N<Kv\).
The shared-gate and coefficient transfer gives
\(s_{\rm circ}\ge\epsilon/(3\gamma)\).
For sufficiently large graphs, take
\(a=c/(2K)\), \(b=c/(3K)\), and absorb the constant three
exactly as in the Theorem~\ref{thm:finite} to obtain \eqref{eq:G5}.

The graph transformation is a fixed local blow-up. Its
polynomial and determinant are uniform whenever the original
graph family is uniform.\end{proof}

The same two-level coefficient and sharp uniform approximation
corollaries hold. This theorem does not assert hardness on
the original \(G\) without the clique blow-up, nor constants
uniform in a gap that shrinks with the graph size.

\section{Finite verification and reproducibility}\label{app:verification}

The proofs above do not depend on experimental extrapolation.
Dedicated computations check the counting conventions, local
gadgets, and circuit transformations on finite examples.
The parity identity \eqref{eq:12} is proved in
Section~\ref{sub:gadget} for every admissible layout and every
bit assignment; the gadget checks below test its implementation
and are not premises of that proof.
Table~\ref{tab:checks} summarizes the principal checks. The accompanying
ancillary files contain the Python implementations, exact result
records, and a manifest of the files used.

\begin{table}[ht]
\centering\small
\begin{tabular}{>{\raggedright\arraybackslash}p{0.33\linewidth}>{\raggedright\arraybackslash}p{0.58\linewidth}}
\toprule
Check & Finite coverage\\
\midrule
BEST counting and sampling &
47,976 local covers and 16,334 last-exit descriptions on
\(C_3,C_4,K_4,K_{3,3}\); exact multiplicities and matrix-tree counts\\
Permutation exposure &
666,600 exact conditional-mean configurations\\
Direct parity gadget &
26,624 bit inputs on 46 layouts; all arcs checked in the physical host\\
Smaller-host implementation &
48 layouts, 12,288 bit inputs, 1,728 locality checks, and
eight exact state-spectrum identities\\
Pair-margin estimate &
2,016 cluster colorings and 93,610 discrete density pairs\\
Balanced shared-gate removal &
2,468 products from 120 DAGs; 6,588 exact output coefficients\\
Shared-gate coefficient charge &
10 global-replacement fixtures, including \(2^{256}\) path multiplicity
and reciprocal \(2^{512}\) rescaling; 20,480 rectangles for
80 signed measures\\
Least-degree component &
12,800 gate-component comparisons in 400 DAGs\\
Counting normalization &
20,789 gate profiles and exact polynomials in 400 DAGs;
10,240 Boolean input points\\
\bottomrule
\end{tabular}
\caption{Finite checks completed for the source arguments.}\label{tab:checks}
\end{table}

The counting check includes a circuit of formal degree
\(4\cdot2^{256}\), whose normalized output has 16 monomials.
The shared-gate check allows multiple productive uses of the selected
gate, and tests the global substitution at all uses.
Incorrect first-exit samplers, root-exit modifications, single-use
charges, and omitted variable-restoration steps were rejected.
The additional charge check distinguishes different DAG nodes that
compute the same polynomial and verifies the level-set decomposition
with exact rational arithmetic. These finite checks supplement
Lemma~\ref{lem:balanced} and Proposition~\ref{prop:charge}.

The toy graphs do not meet the large-degree and spectral hypotheses
of Theorem~\ref{thm:finite}. These computations check finite identities
and implementations, not the universal spectral estimates, concentration
inequalities, or asymptotic circuit lower bound. Reproducibility of a
finite check is also distinct from independent mathematical review.

\begin{thebibliography}{99}
\small
\raggedright

\bibitem{CDM2021}
Arkadev Chattopadhyay, Rajit Datta, and Partha Mukhopadhyay.
Lower bounds for monotone arithmetic circuits via communication complexity.
In \emph{Proceedings of STOC 2021}, pages 786--799, 2021.
\href{https://doi.org/10.1145/3406325.3451069}{doi:10.1145/3406325.3451069}.

\bibitem{CDGM2022}
Arkadev Chattopadhyay, Rajit Datta, Utsab Ghosal, and Partha Mukhopadhyay.
Monotone complexity of spanning tree polynomial re-visited.
In \emph{13th Innovations in Theoretical Computer Science Conference},
volume 215 of \emph{LIPIcs}, pages 39:1--39:21, 2022.
\href{https://doi.org/10.4230/LIPIcs.ITCS.2022.39}{doi:10.4230/LIPIcs.ITCS.2022.39}.

\bibitem{CGM2022}
Arkadev Chattopadhyay, Utsab Ghosal, and Partha Mukhopadhyay.
Robustly separating the arithmetic monotone hierarchy via graph inner-product.
In \emph{42nd IARCS Annual Conference on Foundations of Software
Technology and Theoretical Computer Science},
volume 250 of \emph{LIPIcs}, pages 12:1--12:20, 2022.
\href{https://doi.org/10.4230/LIPIcs.FSTTCS.2022.12}{doi:10.4230/LIPIcs.FSTTCS.2022.12}.

\bibitem{FarrellLevine2015}
Matthew Farrell and Lionel Levine.
Multi-Eulerian tours of directed graphs.
Preprint, 2015.
\href{https://arxiv.org/abs/1509.06237}{arXiv:1509.06237}.

\bibitem{FGK2013}
Sergey Fomin, Dima Grigoriev, and Gleb Koshevoy.
Subtraction-free complexity, cluster transformations, and spanning trees.
Max-Planck-Institut f\"ur Mathematik, Preprint 2013 (36), 2013.
\url{https://webdoc.sub.gwdg.de/ebook/serien/e/mpi_mathematik/2013/36.pdf}.

\bibitem{Hrubes2020}
Pavel Hrube\v{s}.
On \(\epsilon\)-sensitive monotone computations.
\emph{Computational Complexity}, 29, article 6, 2020.
\href{https://doi.org/10.1007/s00037-020-00196-6}{doi:10.1007/s00037-020-00196-6}.

\bibitem{JerrumSnir1982}
Mark Jerrum and Marc Snir.
Some exact complexity results for straight-line computations over semirings.
\emph{Journal of the ACM}, 29(3):874--897, 1982.
\href{https://doi.org/10.1145/322326.322341}{doi:10.1145/322326.322341}.

\bibitem{Jukna2016}
Stasys Jukna.
Lower bounds for monotone counting circuits.
\emph{Discrete Applied Mathematics}, 213:139--152, 2016.
\href{https://doi.org/10.1016/j.dam.2016.04.024}{doi:10.1016/j.dam.2016.04.024}.

\bibitem{LPS1988}
Alexander Lubotzky, Ralph Phillips, and Peter Sarnak.
Ramanujan graphs.
\emph{Combinatorica}, 8(3):261--277, 1988.
\url{https://math.huji.ac.il/~alexlub/PAPERS/ramanujan%20graphs/ramanujanGraphs.pdf}.

\bibitem{MartensMedabalimi2015}
James Martens and Venkatesh Medabalimi.
On the expressive efficiency of sum product networks.
Preprint, version 3, 2015.
\href{https://arxiv.org/abs/1411.7717v3}{arXiv:1411.7717v3}.

\bibitem{Peharz2015}
Robert Peharz, Sebastian Tschiatschek, Franz Pernkopf, and Pedro Domingos.
On theoretical properties of sum-product networks.
In \emph{Proceedings of the Eighteenth International Conference on
Artificial Intelligence and Statistics},
volume 38 of \emph{PMLR}, pages 744--752, 2015.
\url{https://proceedings.mlr.press/v38/peharz15.html}.

\bibitem{ThornerZaman2021}
Jesse Thorner and Asif Zaman.
Refinements to the prime number theorem for arithmetic progressions.
\emph{Mathematische Zeitschrift}, 306(3), article 54, 2024.
\href{https://doi.org/10.1007/s00209-023-03414-3}{doi:10.1007/s00209-023-03414-3}.

\end{thebibliography}


\end{document}
